\documentclass[11pt]{article}
\usepackage{preamble}
\renewcommand{\answer}[1]{}

\begin{document}

\ladderheading{AP Statistics \textperiodcentered\ Unit 1 \textperiodcentered\ Topic 1.6 \textperiodcentered\ B: 2026-09-17 \textperiodcentered\ E: 2026-09-21}{Describing a Distribution}

\focusbanner{TODAY'S PROBLEM}

The histogram shows how long it took 40 students to get to school one morning, in minutes. Two students described the distribution.\par\smallskip \textbf{Ana:} ``Roughly symmetric, centered near 20 minutes, spread from 5 to 45 minutes.''\par \textbf{Ben:} ``Skewed right with a peak in the 10--20 minute interval; most students take under 30 minutes; one student at about 70 minutes is an outlier.''\par
\begin{center}
\begin{tikzpicture}
\begin{axis}[
  width=0.85\linewidth,
  height=2.30in,
  ybar interval, ymin=0, ymax=15,
  xmin=0, xmax=80, xtick={0,10,20,30,40,50,60,70,80},
  xlabel={Commute time (minutes)}, ylabel={Number of students},
  ymajorgrids, tick label style={font=\small}, label style={font=\small},
  bar shift=0pt, x tick label as interval=false
]
\addplot[fill=calloutblue, draw=dayblue] coordinates {
  (0,3)
  (10,13)
  (20,10)
  (30,7)
  (40,4)
  (50,2)
  (60,0)
  (70,1)
  (80,0)
};
\end{axis}
\end{tikzpicture}
\end{center}
\begin{firsttakebox}{FIRST TAKE --- before the video (5 min)}
Whose description do you trust more, Ana's or Ben's? One sentence. No wrong answers yet.
\workspace{1.0in}
\end{firsttakebox}

\begin{callout}{\IconBook\ DESCRIBING A DISTRIBUTION (keep on the page)}{calloutgreen}
Describe \textbf{shape, center, variability, and unusual features} in context, with units. Center describes a typical response; variability describes how spread out values are. A longer right tail means skewed right; a longer left tail means skewed left. An outlier is unusually far from the rest. A gap has no observations. A histogram identifies intervals, not exact individual values.
\end{callout}

\newpage
\textbf{Finish after the video:}\par\smallskip

\dokbadge{1}~\textbf{(a)} Which interval has the tallest bar?\par
\workspace{0.8in}

\dokbadge{2}~\textbf{(b)} In two or three sentences, describe the commute times: shape, a typical value, overall spread, and one unusual feature. Use minutes.\par
\workspace{1.5in}

\dokbadge{3}~\textbf{(c)} Whose description is better supported? Use two graph features to defend your choice. Explain why reporting only a typical commute time would leave a reader with an incomplete picture.\par
\workspace{2.4in}

\begin{sentenceframebox}\raggedright\textbf{Frame:} I trust \blank[0.8in] because the graph shows \blank[1.6in] and \blank[1.6in].\end{sentenceframebox}

\begin{sentenceframebox}\raggedright\textbf{Frame:} A typical time alone would leave out \blank[2in], which matters because \blank[2in].\end{sentenceframebox}

\begin{summaryexitbox}{EXIT --- turn this in}
One thing the video changed about my first take: \hrulefill\par
\workspace{0.4in}
\end{summaryexitbox}

\bankitem{Optional \#1 --- not collected}{\dokbadge{2}~Sketch a dotplot of 15 quiz scores (0--10) that is skewed LEFT, unimodal, and has one low outlier. Then write the one sentence a reader would need in order to describe it in context.}{}

\end{document}
